% =====================================================================
%  1. OBJETIVO
% =====================================================================
\section{Objetivo}
Completar el script \texttt{plantilla\_campos\_potenciales.m} para que el robot navegue de forma reactiva desde un origen hasta un destino sobre el mapa \texttt{mapa1\_150.png}, utilizando \'{u}nicamente las fuerzas de atracci\'{o}n al objetivo y de repulsi\'{o}n de los obst\'{a}culos detectados en cada iteraci\'{o}n. Realizar varios experimentos con diferentes or\'{i}genes y destinos, probar configuraciones de m\'{i}nimo local, y responder a la pregunta del enunciado: \textbf{¿puede el m\'{e}todo encontrar la soluci\'{o}n cambiando solo $\alpha$, $\beta$ o el rango de influencia?}


% =====================================================================
%  2. IMPLEMENTACIÓN
% =====================================================================
\section{Implementaci\'{o}n}

El c\'{o}digo a completar es el bucle de control del robot dentro de la plantilla proporcionada. La estructura de la plantilla (carga del mapa, selecci\'{o}n con \texttt{ginput}, llamada a \texttt{SimulaLidar}, dibujado incremental con \texttt{plot}/\texttt{drawnow}) se mantiene tal cual; solo se a\~{n}aden los c\'{a}lculos de fuerzas y la actualizaci\'{o}n de la pose. Las decisiones no obvias son:

\paragraph{Filtrado de \texttt{NaN} antes de calcular la repulsi\'{o}n.} \texttt{SimulaLidar} (envoltura de \texttt{rayIntersection}) devuelve \texttt{NaN} para los rayos que no impactan dentro del alcance. Si no se descartan con una m\'{a}scara l\'{o}gica antes de calcular distancias, las operaciones vectoriales se contaminan y la repulsi\'{o}n pierde sentido.

\paragraph{Recorte por rango de influencia.} La f\'{o}rmula repulsiva del enunciado se anula exactamente en $d = D$, pero evaluarla para todos los rayos del l\'{a}ser es innecesario. Se filtran previamente los obst\'{a}culos a $d_i < D$ y solo se itera sobre ellos.

\paragraph{Detecci\'{o}n de m\'{i}nimo local por l\'{i}mite de iteraciones.} Se considera que el robot ha quedado atrapado cuando alcanza $1000$ iteraciones sin satisfacer la condici\'{o}n de llegada $\lVert q_\text{dest} - q\rVert < v$. La plantilla ya emite el mensaje correspondiente.


% =====================================================================
%  3. PRUEBAS
% =====================================================================
\section{Pruebas}

Sobre el mapa \texttt{mapa1\_150.png}, con los par\'{a}metros por defecto de la plantilla ($\alpha = 1$, $\beta = 100$, $D = 1.5$, paso $v = 0.4$), se han realizado experimentos con or\'{i}genes y destinos seleccionados con \texttt{ginput}. La Figura~\ref{fig:trayectoria_p3} muestra los dos comportamientos t\'{i}picos.

\begin{figure}[H]
  \centering
  \begin{subfigure}[t]{0.48\linewidth}
    \centering
    \includegraphics[width=\linewidth]{favorable.jpg}
    \caption{Caso favorable: el robot bordea los obst\'{a}culos y alcanza el destino con una trayectoria suave.}
    \label{fig:trayectoria_p3_ok}
  \end{subfigure}\hfill
  \begin{subfigure}[t]{0.48\linewidth}
    \centering
    \includegraphics[width=\linewidth]{minimo_local.jpg}
    \caption{M\'{i}nimo local: atracci\'{o}n y repulsi\'{o}n se cancelan y el robot queda atrapado tras agotar las $1000$ iteraciones.}
    \label{fig:trayectoria_p3_minimo}
  \end{subfigure}
  \caption{Trayectorias reactivas obtenidas con campos potenciales sobre \texttt{mapa1\_150.png}. C\'{i}rculo verde: origen. C\'{i}rculo rojo: destino. Traza roja: recorrido del robot.}
  \label{fig:trayectoria_p3}
\end{figure}

\subsection{Efecto de los par\'{a}metros}
Variando $\alpha$ y $\beta$ se observa el equilibrio entre avance y evitaci\'{o}n: con $\beta$ peque\~{n}o el robot se acerca demasiado a las paredes, con $\beta$ grande la trayectoria se aleja innecesariamente y aparecen oscilaciones. Variando $D$ se controla cu\'{a}ndo empieza a desviarse el robot: $D$ peque\~{n}o produce reacciones bruscas en cercan\'{i}as, $D$ grande adelanta la desviaci\'{o}n. Los valores por defecto ofrecen un compromiso adecuado en este mapa.

\subsection{¿Pueden los par\'{a}metros resolver un m\'{i}nimo local?}
\textbf{No.} Un m\'{i}nimo local es un punto $q^*$ donde $\mathbf{F}_\text{atr}(q^*) + \mathbf{F}_\text{rep}(q^*) = \mathbf{0}$. Cambiar $\alpha$ y $\beta$ \emph{escala} ambas componentes, pero si las fuerzas se cancelan en $q^*$ con una pareja $(\alpha, \beta)$, tambi\'{e}n se cancelan con cualquier $(\lambda\alpha, \lambda\beta)$. Modificar la relaci\'{o}n $\beta/\alpha$ desplaza el punto de cancelaci\'{o}n a otra ubicaci\'{o}n del entorno pero no elimina el m\'{i}nimo, simplemente lo translada. Variar $D$ tampoco rompe la cancelaci\'{o}n: cambia qu\'{e} obst\'{a}culos contribuyen a la repulsi\'{o}n, pero en cualquier configuraci\'{o}n con simetr\'{i}a (cuello en \textsf{U}, paredes enfrentadas) la cancelaci\'{o}n vectorial sigue siendo posible.

Los m\'{i}nimos locales son por tanto una limitaci\'{o}n \textbf{estructural} del m\'{e}todo reactivo puro, no un problema de calibraci\'{o}n. Resolverlos requiere extender el algoritmo (perturbaci\'{o}n aleatoria, campo tangencial, informaci\'{o}n global), v\'{i}a que se explora en Lab-MR 6.


% =====================================================================
%  4. COMENTARIOS FINALES
% =====================================================================
\section{Comentarios finales}

La implementaci\'{o}n cumple lo que pide el enunciado: el robot navega reactivamente, alcanza el destino en configuraciones favorables y queda detectado como atrapado en configuraciones de m\'{i}nimo local. La respuesta a la pregunta del enunciado es negativa: ajustar par\'{a}metros no resuelve los m\'{i}nimos locales, que requieren modificar el algoritmo. Esa extensi\'{o}n se aborda en Lab-MR 6, donde la planificaci\'{o}n global (Dijkstra/A*) act\'{u}a como gu\'{i}a y los campos potenciales se complementan con un campo tangencial y una maniobra de escape por estancamiento.
